Choropleth přeshraničního pracovního dojíždění napříč EU27 ukazuje, že fenomén není omezen na malé otevřené ekonomiky (LU, LI): i~středně velké ekonomiky jako SK a~AT vykazují nadprůměrné hodnoty v~důsledku těsné mzdové proximy k~sousedním trhům práce.
Geografická distribuce potvrzuje tezi o~mzdovém diferenciálu jako hlavním push/pull faktoru: nejintenzivnější přeshraniční dojíždění se koncentruje podél linie nejvyššího mzdového gradientu v~\ac{EU} --- od Pobaltí přes Visegrád až po Balkán, kde jsou diferenciály vůči Německu nebo Skandinávii největší.
Zajímavé je, že ČR jako exportní ekonomika silně integrovaná do německo-rakouských hodnotových řetězců není ve srovnání s~SK výrazněji exponována, což pravděpodobně odráží jazykovou bariéru a~lepší domácí mzdový vývoj v~Praze a~středočeském regionu.
Snižování diferenciálu prostřednictvím silnějšího kolektivního vyjednávání má tedy i~geopolitický rozměr: udržení pracovní síly v~ČR je podmínkou zachování průmyslové výrobní kapacity závislé na zahraniční poptávce.

Choropleth přeshraničního pracovního dojíždění napříč EU27 ukazuje, že fenomén není omezen na malé otevřené ekonomiky (LU, LI): i~středně velké ekonomiky jako SK a~AT vykazují nadprůměrné hodnoty v~důsledku těsné mzdové proximy k~sousedním trhům práce.
Geografická distribuce potvrzuje tezi o~mzdovém diferenciálu jako hlavním push/pull faktoru: nejintenzivnější přeshraniční dojíždění se koncentruje podél linie nejvyššího mzdového gradientu v~\ac{EU} --- od Pobaltí přes Visegrád až po Balkán, kde jsou diferenciály vůči Německu nebo Skandinávii největší.
Zajímavé je, že ČR jako exportní ekonomika silně integrovaná do německo-rakouských hodnotových řetězců není ve srovnání s~SK výrazněji exponována, což pravděpodobně odráží jazykovou bariéru a~lepší domácí mzdový vývoj v~Praze a~středočeském regionu.
Snižování diferenciálu prostřednictvím silnějšího kolektivního vyjednávání má tedy i~geopolitický rozměr: udržení pracovní síly v~ČR je podmínkou zachování průmyslové výrobní kapacity závislé na zahraniční poptávce.

Choropleth přeshraničního pracovního dojíždění napříč EU27 ukazuje, že fenomén není omezen na malé otevřené ekonomiky (LU, LI): i~středně velké ekonomiky jako SK a~AT vykazují nadprůměrné hodnoty v~důsledku těsné mzdové proximy k~sousedním trhům práce.
Geografická distribuce potvrzuje tezi o~mzdovém diferenciálu jako hlavním push/pull faktoru: nejintenzivnější přeshraniční dojíždění se koncentruje podél linie nejvyššího mzdového gradientu v~\ac{EU} --- od Pobaltí přes Visegrád až po Balkán, kde jsou diferenciály vůči Německu nebo Skandinávii největší.
Zajímavé je, že ČR jako exportní ekonomika silně integrovaná do německo-rakouských hodnotových řetězců není ve srovnání s~SK výrazněji exponována, což pravděpodobně odráží jazykovou bariéru a~lepší domácí mzdový vývoj v~Praze a~středočeském regionu.
Snižování diferenciálu prostřednictvím silnějšího kolektivního vyjednávání má tedy i~geopolitický rozměr: udržení pracovní síly v~ČR je podmínkou zachování průmyslové výrobní kapacity závislé na zahraniční poptávce.

Choropleth přeshraničního pracovního dojíždění napříč EU27 ukazuje, že fenomén není omezen na malé otevřené ekonomiky (LU, LI): i~středně velké ekonomiky jako SK a~AT vykazují nadprůměrné hodnoty v~důsledku těsné mzdové proximy k~sousedním trhům práce.
Geografická distribuce potvrzuje tezi o~mzdovém diferenciálu jako hlavním push/pull faktoru: nejintenzivnější přeshraniční dojíždění se koncentruje podél linie nejvyššího mzdového gradientu v~\ac{EU} --- od Pobaltí přes Visegrád až po Balkán, kde jsou diferenciály vůči Německu nebo Skandinávii největší.
Zajímavé je, že ČR jako exportní ekonomika silně integrovaná do německo-rakouských hodnotových řetězců není ve srovnání s~SK výrazněji exponována, což pravděpodobně odráží jazykovou bariéru a~lepší domácí mzdový vývoj v~Praze a~středočeském regionu.
Snižování diferenciálu prostřednictvím silnějšího kolektivního vyjednávání má tedy i~geopolitický rozměr: udržení pracovní síly v~ČR je podmínkou zachování průmyslové výrobní kapacity závislé na zahraniční poptávce.

\input{texparts/python/cross_border_commuting_map}



